\paperchapter{November 2022}{6 November 2022}{A}{50}{16th GSET}
\keytable{\kcell{1}{C} & \kcell{2}{C} & \kcell{3}{D} & \kcell{4}{C} & \kcell{5}{C} & \kcell{6}{B} & \kcell{7}{D} & \kcell{8}{A} & \kcell{9}{B} & \kcell{10}{C} \\
\kcell{11}{D} & \kcell{12}{B} & \kcell{13}{B} & \kcell{14}{C} & \kcell{15}{A} & \kcell{16}{A} & \kcell{17}{D} & \kcell{18}{C} & \kcell{19}{B} & \kcell{20}{C} \\
\kcell{21}{Z} & \kcell{22}{D} & \kcell{23}{A} & \kcell{24}{C} & \kcell{25}{A} & \kcell{26}{B} & \kcell{27}{C} & \kcell{28}{A} & \kcell{29}{A} & \kcell{30}{D} \\
\kcell{31}{B} & \kcell{32}{A} & \kcell{33}{B} & \kcell{34}{D} & \kcell{35}{B} & \kcell{36}{D} & \kcell{37}{B} & \kcell{38}{B} & \kcell{39}{D} & \kcell{40}{B} \\
\kcell{41}{A} & \kcell{42}{D} & \kcell{43}{B} & \kcell{44}{D} & \kcell{45}{C} & \kcell{46}{D} & \kcell{47}{C} & \kcell{48}{C} & \kcell{49}{D} & \kcell{50}{D}}
\begin{context}
\textbf{Directions for Q.~1 to 5:} Read the following passage and answer the questions.\par\smallskip
At first sight silence would seem to be the essence of non communication. Yet we know that in some circumstances silence can speak volumes. It is also a potent tool in the communication tool box, for the speaker, as it adds emphasis to point, allowing time for its full significance to sink in. Handled skillfully, it can build up anticipation for what is to come the pregnant pause. It gives opportunity to make eye contact with audience and gauge their reactions. It is an essential component of the timing of one's delivery. For the listener silence can buy time, either to digest what has been said or to consider a response, or both. Silence will, often lure the speaker into saying more than he or she may have intended.
\end{context}
\begin{question}{1}
{Silence is an essential component of\par}
\optstwo{Non-communication}{Gauging reaction of people}{The timing of one's delivery}{Distraction}
\answer{C}{The timing of one's delivery}
\begin{explain}
The passage states directly: ``It is an essential component of \textbf{the timing of one's delivery}.''
\why{``Non-communication'' is what silence \emph{seems} to be at first sight, a view the passage then rejects. Gauging reactions is an \emph{opportunity} that silence gives. Distraction is not mentioned.}
\end{explain}
\end{question}

\begin{question}{2}
{According to the author, the pregnant pause\par}
\optstwo{Breaks up communication}{Helps communication}{Build up anticipation}{Leads to confusion}
\answer{C}{Build up anticipation}
\begin{explain}
``Handled skillfully, it can \textbf{build up anticipation} for what is to come: the pregnant pause.'' A pregnant pause is a deliberate silence full of meaning that makes the audience wait for what follows.
\end{explain}
\end{question}

\begin{question}{3}
{For the listener silence can\par}
\optstwo{Emphasize a point}{Be a potent tool}{Speak volumes}{Buy time}
\answer{D}{Buy time}
\begin{explain}
``For the listener silence can \textbf{buy time}, either to digest what has been said or to consider a response.''
\why{Emphasising a point, being a potent tool and speaking volumes are described as uses of silence for the \emph{speaker} or in general.}
\end{explain}
\end{question}

\begin{question}{4}
{Silence seems to be essence of\par}
\optstwo{Communication}{Transmission}{Non-Communication}{Communication tool box}
\answer{C}{Non-Communication}
\begin{explain}
``At first sight silence would seem to be the essence of \textbf{non communication}.''
\end{explain}
\end{question}

\begin{question}{5}
{What provides an opportunity to make eye contact with the audience ?\par}
\optsfour{Communication}{Discussion}{Silence}{Gesture}
\answer{C}{Silence}
\begin{explain}
``It [silence] gives opportunity to make \textbf{eye contact} with audience and gauge their reactions.''
\end{explain}
\end{question}

\begin{question}{6}
{Which of the following are considered as both input and output devices ?\newline a. Touch screen\newline b. USB memory devices\newline c. CD drive\newline d. Screen projector\par}
\optsfour{a, c and d}{a, b and c}{a, b and d}{b, c and d}
\answer{B}{a, b and c}
\begin{explain}
Devices that both \emph{accept} and \emph{deliver} data are input/output devices:
\begin{itemize}
\item \textbf{Touch screen}: displays output and takes touch input.
\item \textbf{USB memory (pen drive)}: data can be written to it and read from it.
\item \textbf{CD drive} (read/write): reads from and writes to discs.
\end{itemize}
A \textbf{screen projector} only displays output. Hence a, b and c.
\end{explain}
\end{question}

\begin{question}{7}
{Which of the following statements are true ?\newline a. Cache memory is used to enhance the performance of computer system.\newline b. Primary memory is also known as main memory.\newline c. Hard disk is an example of secondary memory.\newline d. Static Random Access Memory (SRAM) is not a Cache memory.\par}
\optsone{a and b are true, while c and d are false}{a and c are true, while b and d are false}{b, c and d are true, while a is false}{a, b and c are true, while d is false}
\answer{D}{a, b and c are true, while d is false}
\begin{explain}
(a) True: cache is a small, very fast memory between the CPU and RAM that speeds up processing.
(b) True: primary memory (RAM/ROM) is also called main memory.
(c) True: the hard disk is secondary (auxiliary) storage.
(d) False: \textbf{SRAM is precisely the memory used for cache}, because it is fast and needs no refreshing. DRAM is used for main memory.
Hence a, b and c are true and d is false.
\end{explain}
\end{question}

\begin{question}{8}
{Which of the following are characteristics of CC-BY SA?\newline a. Allows to copy and publish\newline b. Attribution required\newline c. Cannot be used commercially\newline d. Cannot be modified and adapted\newline e. Cannot change license\par}
\optsfour{a, b and e}{a, b, d and e}{a, c, d and e}{b, d and e}
\answer{A}{a, b and e}
\begin{explain}
\textbf{CC BY-SA} (Creative Commons \emph{Attribution--ShareAlike}):
\begin{itemize}
\item allows copying, redistribution and publishing (a);
\item requires \textbf{attribution} to the creator (b);
\item permits commercial use, so (c) is wrong;
\item permits modification and adaptation, so (d) is wrong;
\item derivatives must be released \textbf{under the same licence}, so the licence cannot be changed (e).
\end{itemize}
The answer is a, b and e.
\tip{CC elements: BY = attribution, SA = share-alike, NC = non-commercial, ND = no derivatives. CC0 = public-domain dedication.}
\end{explain}
\end{question}

\begin{question}{9}
{In relation to NMEICT mission document, which of the following is not true for opportunities of ICT in education?\par}
\optsone{Availability of EduSAT}{High power consumption connectivity and computing devices}{The abundance of knowledge on the internet}{Rapidly expanding network of cyber kiosks and cyber cafes}
\answer{B}{High power consumption connectivity and computing devices}
\begin{explain}
The \emph{National Mission on Education through ICT (NMEICT)} lists as opportunities the availability of EduSAT, the abundance of knowledge on the internet, and the fast-growing network of cyber kiosks and cafes. The mission document looks for \textbf{low}-power, low-cost connectivity and devices (e.g.\ the low-cost \emph{Aakash} tablet), so ``high power consumption'' is \emph{not} an opportunity.
\end{explain}
\end{question}

\begin{question}{10}
{Which of the following has been conceived as 'One Stop Education Portal'?\par}
\optsfour{SAMAGRA}{SAMARTH}{SAKSHAT}{DIKSHA}
\answer{C}{SAKSHAT}
\begin{explain}
\textbf{SAKSHAT}, launched on 30 October 2006 by MHRD, was conceived as a \textbf{``One Stop Education Portal''} addressing the educational needs of learners from school to postgraduate level. It is the portal of the NMEICT.
\why{DIKSHA is the national platform for school teachers and students (2017); SAMARTH is an e-governance ERP for universities; Samagra Shiksha is a scheme for school education.}
\end{explain}
\end{question}

\begin{question}{11}
{Look this series 40, 38, 34, 32, 28, what number should come next?\par}
\optsfour{24}{22}{27}{26}
\answer{D}{26}
\begin{explain}
The differences alternate between $-2$ and $-4$:
\[40\xrightarrow{-2}38\xrightarrow{-4}34\xrightarrow{-2}32\xrightarrow{-4}28\xrightarrow{-2}\mathbf{26}\]
\end{explain}
\end{question}

\begin{question}{12}
{Find the next term of given series.\newline ABD, DGK, HMS, MTB, SBL, \blank{}\par}
\optsfour{ZKU}{ZKW}{XKW}{ZAB}
\answer{B}{ZKW}
\begin{explain}
Treat the three letter positions as three separate series (A=1, \ldots, Z=26, counting cyclically):
\begin{itemize}
\item 1st letters: A, D, H, M, S (+3, +4, +5, +6), so next $+7$ gives \textbf{Z}.
\item 2nd letters: B, G, M, T, B (+5, +6, +7, +8), so next $+9$: B(2)$+9=$ \textbf{K}.
\item 3rd letters: D, K, S, B, L (+7, +8, +9, +10), so next $+11$: L(12)$+11=$ \textbf{W}.
\end{itemize}
The next term is \textbf{ZKW}.
\end{explain}
\end{question}

\begin{question}{13}
{Find the next term of given series.\newline C4X, F9U, I16R, \blank{}\par}
\optsfour{L25P}{L25O}{L27P}{L29S}
\answer{B}{L25O}
\begin{explain}
\begin{itemize}
\item First letters: C, F, I, \textbf{L} (+3).
\item Numbers: $4=2^2$, $9=3^2$, $16=4^2$, so next $5^2=$ \textbf{25}.
\item Last letters: X, U, R, \textbf{O} ($-3$).
\end{itemize}
The next term is \textbf{L25O}.
\end{explain}
\end{question}

\begin{question}{14}
{Which of the following is in descending order ?\par}
\optstwo{5/8, 9/13, 11/17}{5/8, 11/17, 9/13}{9/13, 11/17, 5/8}{11/17, 9/13, 5/8}
\answer{C}{9/13, 11/17, 5/8}
\begin{explain}
Decimal values: $\tfrac58=0.625$, $\tfrac{9}{13}\approx0.692$, $\tfrac{11}{17}\approx0.647$.
Descending order: $\tfrac{9}{13}>\tfrac{11}{17}>\tfrac58$, which is option (C).
\end{explain}
\end{question}

\begin{question}{15}
{4/7 of 2/3 of 5/6 of 5/8 of 1008 is\par}
\optsfour{200}{144}{64}{400}
\answer{A}{200}
\begin{explain}
\[\frac47\times\frac23\times\frac56\times\frac58\times1008=\frac{4\times2\times5\times5}{7\times3\times6\times8}\times1008=\frac{200}{1008}\times1008=\mathbf{200}\]
\end{explain}
\end{question}

\begin{question}{16}
{Which of the following pollutants pose the highest risk to human health and environment ?\par}
\optstwo{Particulate matter}{Nitrogen dioxide}{Sulphur dioxide}{Carbon monoxide}
\answer{A}{Particulate matter}
\begin{explain}
\textbf{Particulate matter}, especially PM$_{2.5}$, is regarded by WHO as the pollutant with the greatest health impact. The fine particles penetrate deep into the lungs and bloodstream and cause respiratory and cardiovascular disease and premature death.
\end{explain}
\end{question}

\begin{question}{17}
{Which of the following is not a priority area for Education for Sustainable Development' for 2030?\par}
\optstwo{Advancing Policy}{Transforming Learning Environments}{Empowering and mobilizing youth}{Building capacities of administrators}
\answer{D}{Building capacities of administrators}
\begin{explain}
UNESCO's \textbf{ESD for 2030} framework has five priority action areas: (1) advancing policy, (2) transforming learning environments, (3) building capacities of \emph{educators}, (4) empowering and mobilising youth, and (5) accelerating local-level actions. ``Building capacities of \textbf{administrators}'' is not one of them.
\end{explain}
\end{question}

\begin{question}{18}
{Match the following and choose the correct option given below.\\[2pt]
\begin{tabular}{@{}ll@{}}
\textbf{Movement} & \textbf{Activist}\\
I. Chipko Movement & a. M.~S.~Swaminathan\\
II. Narmada Bachao Movement & b. Medha Patkar\\
III. Green Revolution & c. Amrita Devi\\
IV. Bishnoi Movement & d. Sunderlal Bahuguna\\
\end{tabular}\par}
\optstwo{I-b, II-a, III-d, IV-c}{I-c, II-b, III-a, IV-d}{I-d, II-b, III-a, IV-c}{I-a, II-b, III-c, IV-d}
\answer{C}{I-d, II-b, III-a, IV-c}
\begin{explain}
\begin{itemize}
\item \textbf{Chipko} (1973, Uttarakhand): Sunderlal Bahuguna (also Chandi Prasad Bhatt and Gaura Devi), so I-d.
\item \textbf{Narmada Bachao Andolan} (1985): Medha Patkar, so II-b.
\item \textbf{Green Revolution} (1960s): M.~S.~Swaminathan, so III-a.
\item \textbf{Bishnoi movement} (1730, Khejarli, Rajasthan): Amrita Devi Bishnoi and 363 Bishnois gave their lives to save Khejri trees, so IV-c.
\end{itemize}
\end{explain}
\end{question}

\begin{question}{19}
{"The Future We Want" is an outcome document of\par}
\optsone{Brundtland Commission}{United Nations Conference on Sustainable Development}{World Commission on Environment and Development}{United Nations Environment Programme}
\answer{B}{United Nations Conference on Sustainable Development}
\begin{explain}
\emph{``The Future We Want''} is the outcome document of the \textbf{United Nations Conference on Sustainable Development (Rio+20)}, held in Rio de Janeiro in June 2012. It launched the process that led to the SDGs.
\why{The Brundtland Commission (World Commission on Environment and Development) produced \emph{Our Common Future} (1987).}
\end{explain}
\end{question}

\begin{question}{20}
{The Nagoya Protocol on 'Access to Genetic Resources and the fair and equitable sharing of Benefits' is an\par}
\optsone{International agreement that aims to lessen water pollution}{International agreement which aims at sharing the benefits arising from the utilization of genetic resources in all possible ways}{International agreement which aims at sharing the benefits arising from the utilization of genetic resources in a fair and equitable way}{International agreement which aims to ensure the safe handling, transport and use of living modified organisms (LMOs)}
\answer{C}{International agreement which aims at sharing the benefits arising from the utilization of genetic resources in a fair and equitable way}
\begin{explain}
The \textbf{Nagoya Protocol} (adopted 2010, in force 2014), a supplementary agreement to the Convention on Biological Diversity, concerns \textbf{Access and Benefit-Sharing (ABS)}: the benefits arising from the use of genetic resources must be shared in a \textbf{fair and equitable} way.
\why{Option (D) describes the \emph{Cartagena Protocol on Biosafety} (2000), which covers living modified organisms.}
\end{explain}
\end{question}

\begin{question}{21}
{Study the chart and answer the question. The pie chart given here represents the domestic expenditure of a family in percent. If the total monthly income of the family is Rs.~55,660/-, after PF deductions and payment of the house rent, the total monthly income of the family remains\par}
\qfigure{\begin{tikzpicture}[scale=0.9]
\fill[brand!15] (0,0) -- (160:1.4) arc (160:235.6:1.4) -- cycle;
\draw (0,0) circle (1.4);
\draw (0,0) -- (160:1.4); \draw (0,0) -- (235.6:1.4);
\node[font=\scriptsize,align=center] at (198:0.85) {Child\\Educ.\\21\%};
\node[font=\scriptsize,align=left,right] at (1.7,0) {Other slices: PF, Food,\\House Rent, Entertainment,\\Medical (values not legible\\in the question booklet)};
\end{tikzpicture}}
\optsfour{38662}{38862}{38962}{38752}
\dropped{}
\begin{explain}
Only the \emph{Child Education = 21\%} slice was legibly printed. The PF and house-rent percentages, which are needed to answer, could not be read from the chart, so the question could not be solved and the agency \textbf{dropped} it.
For reference, the options are about 70\% of the income: $0.70\times55{,}660=38{,}962$ (option C). This would be the answer if PF and house rent together formed 30\% of expenditure.
\end{explain}
\end{question}

\begin{question}{22}
{Choose the graph that best fits with this statement: "House prices rose in January and February, fell slightly in March and are now holding steady".\par}
\optsfour{\tikz[baseline=3mm,scale=0.45]{\draw[->](0,0)--(2.6,0);\draw[->](0,0)--(0,1.6);\draw[thick,brand](0,1)--(1.8,1);}}{\tikz[baseline=3mm,scale=0.45]{\draw[->](0,0)--(2.6,0);\draw[->](0,0)--(0,1.6);\draw[thick,brand](0,0.8) .. controls (1.2,0.8) and (1.8,1) .. (2,1.5);}}{\tikz[baseline=3mm,scale=0.45]{\draw[->](0,0)--(2.6,0);\draw[->](0,0)--(0,1.6);\draw[thick,brand](0,0)--(1.9,1.5);}}{\tikz[baseline=3mm,scale=0.45]{\draw[->](0,0)--(2.6,0);\draw[->](0,0)--(0,1.6);\draw[thick,brand](0,0.6) .. controls (0.4,1.3) and (0.6,1.4) .. (0.9,0.9) .. controls (1.1,0.6) and (1.3,0.9) .. (1.6,0.95) -- (2.3,0.95);}}
\answer{D}{\tikz[baseline=3mm,scale=0.45]{\draw[->](0,0)--(2.6,0);\draw[->](0,0)--(0,1.6);\draw[thick,brand](0,0.6) .. controls (0.4,1.3) and (0.6,1.4) .. (0.9,0.9) .. controls (1.1,0.6) and (1.3,0.9) .. (1.6,0.95) -- (2.3,0.95);}}
\begin{explain}
The statement has three phases. Prices \textbf{rise} (January--February), then \textbf{fall slightly} (March), then stay \textbf{steady} (now). Only graph (D) shows a rise, a small dip and then a flat line.
\why{(A) is flat throughout; (B) and (C) keep rising, with (B) accelerating and (C) at a constant rate.}
\end{explain}
\end{question}

\begin{question}{23}
{A researcher uses systematic random sampling from a company directory that has 500 employees listed in alphabetical order. If the desired samples size is 20, and the first name to be selected is number 9, which of the following will not be selected ?\par}
\optsfour{89}{134}{309}{484}
\answer{A}{89}
\begin{explain}
Sampling interval $k=\dfrac{N}{n}=\dfrac{500}{20}=25$. Selected numbers: $9, 34, 59, 84, 109, 134, \ldots, 309, \ldots, 484$, i.e.\ $9+25m$.
\begin{itemize}
\item $134=9+25\times5$ \checkmark, $309=9+25\times12$ \checkmark, $484=9+25\times19$ \checkmark
\item $89-9=80$ is not a multiple of 25, so \textbf{89 is not selected}.
\end{itemize}
\end{explain}
\end{question}

\begin{question}{24}
{The bar graph given below shows the Foreign Exchange Reserves of a Country (in millions US\$) from 1991-92 to 1998-99. The ratio of the number of years in which the foreign exchange reserves are above the average reserves, to those in which the reserves are below the average reserves is\par}
\qfigure{\begin{tikzpicture}
\begin{axis}[ybar,width=10cm,height=4.4cm,bar width=11pt,ymin=0,ymax=5800,
  symbolic x coords={91-92,92-93,93-94,94-95,95-96,96-97,97-98,98-99},xtick=data,
  nodes near coords,every node near coord/.append style={font=\tiny},
  tick label style={font=\scriptsize},ylabel={\scriptsize million US\$},
  title={\scriptsize Foreign Exchange Reserve of a Country},axis x line*=bottom,axis y line*=left]
\addplot[fill=brand!70,draw=brand] coordinates {(91-92,2640) (92-93,3720) (93-94,2520) (94-95,3360) (95-96,3120) (96-97,4320) (97-98,5040) (98-99,3120)};
\draw[accent,dashed,thick] (axis cs:91-92,3480) -- (axis cs:98-99,3480);
\end{axis}
\end{tikzpicture}}
\optsfour{2:6}{3:4}{3:5}{4:4}
\answer{C}{3:5}
\begin{explain}
\[\text{Average}=\frac{2640+3720+2520+3360+3120+4320+5040+3120}{8}=\frac{27840}{8}=3480\]
Above the average (dashed line): 3720, 4320 and 5040, which is \textbf{3 years}. Below: 2640, 2520, 3360, 3120 and 3120, which is \textbf{5 years}. The ratio is \textbf{3 : 5}.
\end{explain}
\end{question}

\begin{question}{25}
{What is the major purpose of Data Interpretation?\par}
\optsone{Getting meaning from the data and arriving as a conclusion}{Classification of data}{Analysis of data using suitable statistical tool}{Comparing different classified data}
\answer{A}{Getting meaning from the data and arriving as a conclusion}
\begin{explain}
\textbf{Data interpretation} means drawing meaning from analysed data, relating it to the research questions or hypotheses and \textbf{arriving at conclusions}. Classification, statistical analysis and comparison are \emph{earlier} steps that interpretation builds upon.
\end{explain}
\end{question}

\begin{question}{26}
{The main objective of formative evaluation\par}
\optsone{To find out the difficulties in learning}{To enhance students' learning}{To promote students to the next class}{To develop cooperation among students}
\answer{B}{To enhance students' learning}
\begin{explain}
\textbf{Formative evaluation} (Michael Scriven, 1967) is carried out \emph{during} instruction to give continuous feedback, so that teaching can be adjusted and \textbf{students' learning improved}. It is assessment \emph{for} learning.
\why{Identifying learning difficulties is \emph{diagnostic} evaluation; promotion to the next class is based on \emph{summative} evaluation.}
\end{explain}
\end{question}

\begin{question}{27}
{Which of the following options is not befitting to online teaching concept ?\par}
\optstwo{Swayamprabha}{Massive Open Online Courses}{Shodh Ganga}{Google meet}
\answer{C}{Shodh Ganga}
\begin{explain}
\textbf{Shodhganga} (INFLIBNET) is a digital \emph{repository of Indian PhD theses}; it is not an online teaching platform.
\why{SWAYAM Prabha (educational DTH), MOOCs and Google Meet (live classes) are all used for online or distance teaching.}
\end{explain}
\end{question}

\begin{question}{28}
{Which of the following statements about individual differences are true ?\newline 1. Every child is unique and different from each other.\newline 2. The same level of intelligence, the achievement level of student may be different.\newline 3. Heredity and environment both are the reason for individual differences.\newline 4. Individual differences are related to psychology only.\newline 5. Individual differences are there but it does not make any difference.\par}
\optsfour{1, 2 and 3}{1, 2 and 4}{2, 3 and 5}{3, 4 and 5}
\answer{A}{1, 2 and 3}
\begin{explain}
True statements:
\begin{itemize}
\item every child is unique (1);
\item students of the same intelligence may achieve differently, because of motivation, study habits and so on (2);
\item heredity and environment together cause individual differences (3).
\end{itemize}
Statement 4 is wrong because differences are physical, social and emotional as well as psychological. Statement 5 is wrong because differences do affect learning and teaching. The answer is 1, 2 and 3.
\end{explain}
\end{question}

\begin{question}{29}
{One of your student remains quite in the class discussion while he/she is very talkative and playful outside the class, what would you do?\newline 1. Try to find why he/she remains quite in the class discussion and talk to him/her.\newline 2. Assign him/her some topic to speak in the class.\newline 3. You will ask question to the student so he/she has to speak.\newline 4. Let him/her be quite as it is his/her nature.\newline 5. Send him/her to playground and let him play where he/she is comfortable.\par}
\optsfour{1, 2 and 3}{2, 3 and 4}{1, 3 and 5}{2, 4 and 5}
\answer{A}{1, 2 and 3}
\begin{explain}
A good teacher (i) finds out \emph{why} the student is quiet in class and talks to him or her, (ii) gives the student a topic to present so as to build confidence, and (iii) involves the student by asking questions. Leaving the student alone (4) or sending the student to play (5) neglects the teacher's responsibility. The answer is 1, 2 and 3.
\end{explain}
\end{question}

\begin{question}{30}
{Bright student causes disturbance in showing of his/her knowledge while you are teaching, what would you do ?\newline 1. Give him/her some work that he can not do so he realise that he/she doesn't know everything.\newline 2. Give him/her some difficult work in the area of his/her interest.\newline 3. Give him/her some books to read in the area of his/her interest.\newline 4. You will connect his/her points in your talk.\newline 5. Complain against him/her to his/her parents.\par}
\optsfour{1, 2 and 4}{2, 4 and 5}{2, 3 and 5}{2, 3 and 4}
\answer{D}{2, 3 and 4}
\begin{explain}
Constructive handling of a bright student who disrupts class: set challenging work in the area of interest (2), suggest enrichment reading (3), and connect the student's points with the lesson so the contribution is valued (4).
Setting impossible work to humiliate the student (1) or complaining to the parents (5) is negative. The answer is 2, 3 and 4.
\end{explain}
\end{question}

\begin{question}{31}
{Deep Ecology has focused upon\newline a. Ecosophy\newline b. Sustainable life style\newline c. Richness and diversity of life\newline d. Deforestation\newline e. Only abiotic ecosystem\newline Choose the correct option given below.\par}
\optstwo{Options b, d and e are correct}{Options a, c and b are correct}{Options a, b and e are correct}{Options b, c and e are correct}
\answer{B}{Options a, c and b are correct}
\begin{explain}
\textbf{Deep ecology} (Arne Naess, 1973) holds that all living beings have intrinsic value. It emphasises \textbf{ecosophy} (ecological wisdom or harmony), \textbf{richness and diversity of life}, and a \textbf{simple, sustainable lifestyle}. It is not limited to the abiotic ecosystem. The answer is a, b and c.
\end{explain}
\end{question}

\begin{question}{32}
{As per the latest report on 'All India Survey on Higher Education (2019-2020)', the Gender Parity Index in India is\par}
\optsfour{1.01}{1.02}{1.11}{1.21}
\answer{A}{1.01}
\begin{explain}
According to \textbf{AISHE 2019--20}, the \textbf{Gender Parity Index (GPI)} in higher education was \textbf{1.01}. This was the first time that female GER exceeded male GER. Total GER was 27.1\%.
\end{explain}
\end{question}

\begin{question}{33}
{Higher Education 4.0 has emphasized more upon\par}
\optstwo{Synchronous learning}{Blended mode of learning}{Face to face learning}{Adult and continuing learning}
\answer{B}{Blended mode of learning}
\begin{explain}
\emph{Higher Education 4.0}, aligned with Industry 4.0, promotes technology-integrated, flexible, personalised learning. It stresses \textbf{blended learning}, combining online and face-to-face modes.
\end{explain}
\end{question}

\begin{question}{34}
{Match the following given in List--I and List--II and select the correct option given below.\\[2pt]
\begin{tabular}{@{}p{0.3\linewidth}p{0.62\linewidth}@{}}
\textbf{List--I} & \textbf{List--II}\\
1. National Digital Library & a. Plenty of international electronic journals and e-books are made available to all HEIs\\
2. National Academic Depository & b. To incorporate Robotics in Engineering Education\\
3. e-Yantra & c. To facilitate digital issuance, storage, access and verification of Academic credentials\\
4. e-Shodh Sindhu & d. To develop a framework of virtual repository of learning resources with a single window search facility\\
\end{tabular}\par}
\optstwo{1-d, 2-b, 3-a, 4-c}{1-d, 2-a, 3-c, 4-b}{1-b, 2-c, 3-d, 4-a}{1-d, 2-c, 3-b, 4-a}
\answer{D}{1-d, 2-c, 3-b, 4-a}
\begin{explain}
\begin{itemize}
\item \textbf{National Digital Library of India} (IIT Kharagpur): a virtual repository of learning resources with single-window search, so 1-d.
\item \textbf{National Academic Depository (NAD)}: digital issue, storage, access and verification of academic awards, so 2-c.
\item \textbf{e-Yantra} (IIT Bombay): brings robotics into engineering education, so 3-b.
\item \textbf{e-ShodhSindhu} (INFLIBNET): e-journals and e-books for HEIs, so 4-a.
\end{itemize}
\end{explain}
\end{question}

\begin{question}{35}
{Which of the following core principle has not been mentioned in National Education Policy 2020?\par}
\optstwo{Extensive use of Technology}{Respect for Global Diversity}{Synergy in curriculum}{Creativity and critical thinking}
\answer{B}{Respect for Global Diversity}
\begin{explain}
The fundamental principles in NEP~2020 include extensive use of technology, creativity and critical thinking, multidisciplinarity and a holistic approach (synergy in curriculum), and ``respect for diversity and respect for the local context''. \textbf{``Respect for global diversity''} is not among the stated principles.
\end{explain}
\end{question}

\begin{question}{36}
{Match Section--I with Section--II.\\[2pt]
\begin{tabular}{@{}ll@{}}
\textbf{Section--I} & \textbf{Section--II}\\
1. Haptics & a. Use of body movement\\
2. Proxemics & b. Use of space and effects\\
3. Paralanguage & c. Communication through touching\\
4. Kinesics & d. Use of pitch and volume of voice\\
\end{tabular}\par}
\optstwo{1-a, 2-b, 3-c, 4-d}{1-b, 2-a, 3-d, 4-c}{1-c, 2-d, 3-a, 4-b}{1-c, 2-b, 3-d, 4-a}
\answer{D}{1-c, 2-b, 3-d, 4-a}
\begin{explain}
Haptics is touch (1-c); Proxemics is space (2-b); Paralanguage covers pitch, volume, tone and rate of voice (3-d); Kinesics is body movement (4-a). The match is 1-c, 2-b, 3-d, 4-a.
\end{explain}
\end{question}

\begin{question}{37}
{Which of the following are the example of communicative tool?\newline 1. Teleconferencing\newline 2. Voice mail\newline 3. Chat\newline 4. Multimedia Encyclopedia\par}
\optsfour{1, 3 and 4}{1, 2 and 3}{1, 2 and 4}{2, 3 and 4}
\answer{B}{1, 2 and 3}
\begin{explain}
Teleconferencing, voice mail and chat are \textbf{communication tools}: they carry messages between people, synchronously or asynchronously. A \textbf{multimedia encyclopaedia} is an information or reference resource, not a communication tool. The answer is 1, 2 and 3.
\end{explain}
\end{question}

\begin{question}{38}
{Which is not a characteristics of Effective Communication?\par}
\optstwo{Completeness}{Cultural heterogeneity}{Concreteness}{Correctness}
\answer{B}{Cultural heterogeneity}
\begin{explain}
The \textbf{7 Cs of effective communication} are Clarity, Conciseness, Concreteness, Correctness, Coherence, Completeness and Courtesy (some versions add Consideration). \textbf{Cultural heterogeneity} is not one of them; cultural differences are in fact a potential \emph{barrier}.
\end{explain}
\end{question}

\begin{question}{39}
{Which of the following is not a component of a communication system ?\par}
\optsfour{Sender}{Message}{Feedback}{Reinforcement}
\answer{D}{Reinforcement}
\begin{explain}
The components of a communication system are \textbf{sender (source), message, encoding, channel, receiver, decoding and feedback}, along with noise and context. \textbf{Reinforcement} is a learning-theory concept (Skinner), not a component of communication.
\end{explain}
\end{question}

\begin{question}{40}
{What are the causes of semantic barriers in communication ?\newline 1. Poorly expressed message\newline 2. Use of technical language\newline 3. Unnatural assumption\newline 4. Lack of attention\par}
\optsfour{2, 3, 4}{1, 2, 3}{1, 3, 4}{1, 2, 4}
\answer{B}{1, 2, 3}
\begin{explain}
\textbf{Semantic barriers} are barriers of \emph{meaning}. They include a badly expressed message (1), technical language or jargon (2), and unclarified or wrong assumptions (3). \textbf{Lack of attention} (4) is a \emph{psychological} barrier. The answer is 1, 2 and 3.
\end{explain}
\end{question}

\begin{question}{41}
{Given below are two statements - One is labeled as Assertion (A) and the other is labeled as Reason (R).\newline Assertion (A): Government of India has established higher educational bodies to ensure the quality of education.\newline Reason (R): There are some educational institutions in India which do not provide quality education.\newline In the context of the above two statements, choose the correct option.\par}
\optsone{Both (A) and (R) are true and (R) is the correct explanation of (A)}{Both (A) and (R) are true but (R) is not the correct explanation of (A)}{(A) is true, but (R) is false}{(A) is false, but (R) is true}
\answer{A}{Both (A) and (R) are true and (R) is the correct explanation of (A)}
\begin{explain}
The Government has set up regulatory and quality-assurance bodies (UGC, AICTE, NAAC, NBA, NCTE and others) because some institutions fall short on quality. The Reason gives exactly why these bodies were established, so both are true and R explains A.
\end{explain}
\end{question}

\begin{question}{42}
{Which of the following statements are true in the context of Vyapti?\newline 1. It is the relation between middle and major term in the logical ground of inference.\newline 2. It means a state of pervasion.\newline 3. Vyapti serves a basis of Nyaya syllogism.\par}
\optsfour{1 and 2}{2 and 3}{1 and 3}{1, 2 and 3}
\answer{D}{1, 2 and 3}
\begin{explain}
In Nyaya logic, \textbf{Vyapti} (invariable concomitance) is:
\begin{itemize}
\item the relation between the \emph{middle term} (hetu, e.g.\ smoke) and the \emph{major term} (sadhya, e.g.\ fire) (1);
\item literally a state of \emph{pervasion}, in which smoke is pervaded by fire (2);
\item the logical basis of inference and hence of the Nyaya syllogism, expressed in its \emph{Udaharana} step (3).
\end{itemize}
All three statements are true.
\end{explain}
\end{question}

\begin{question}{43}
{Find the next term of given series.\newline ATTRIBUTION, TTRIBUTIO, RIBUTIO, IBUTI.\ldots{}?\par}
\optsfour{BUT}{UTI}{UT}{IBU}
\answer{B}{UTI}
\begin{explain}
Letters are removed alternately from both ends and from the front:
\begin{center}
ATTRIBUTION $\xrightarrow{\text{drop 1st \& last}}$ TTRIBUTIO $\xrightarrow{\text{drop first 2}}$ RIBUTIO $\xrightarrow{\text{drop 1st \& last}}$ IBUTI $\xrightarrow{\text{drop first 2}}$ \textbf{UTI}
\end{center}
\end{explain}
\end{question}

\begin{question}{44}
{In a certain relationship\newline $S\times T$ means S is brother of T\newline S-T means S is mother of T\newline $S+T$ means S is father of T\newline Answer the following question based on data above.\newline In the pattern $M+N+O$ how is M related to O?\par}
\optsfour{Uncle}{Father}{Cousin}{Grandfather}
\answer{D}{Grandfather}
\begin{explain}
$M+N$: M is the father of N. $N+O$: N is the father of O. So M is the father of O's father, i.e.\ O's \textbf{grandfather}.
\end{explain}
\end{question}

\begin{question}{45}
{Which of the following diagrams indicates the best relation between Travelers, Train and Bus?\par}
\optsfour{\tikz[baseline=-0.5ex,scale=0.3]{\draw (0,0) circle(1.2); \draw (0,0) circle(0.6); \draw (2.2,0) ellipse(0.6 and 0.4);}}{\tikz[baseline=-0.5ex,scale=0.3]{\draw (0,0.5) circle(0.7); \draw (1.3,0.5) circle(0.7); \draw (0.65,-0.6) circle(0.7);}}{\tikz[baseline=-0.5ex,scale=0.3]{\draw (0,0) circle(1.3); \draw (-0.55,0) circle(0.5); \draw (0.55,0) circle(0.5);}}{\tikz[baseline=-0.5ex,scale=0.3]{\draw (0,0) circle(0.8); \draw (1.1,0) circle(0.8); \draw (2.2,0) circle(0.8);}}
\answer{C}{\tikz[baseline=-0.5ex,scale=0.3]{\draw (0,0) circle(1.3); \draw (-0.55,0) circle(0.5); \draw (0.55,0) circle(0.5);}}
\begin{explain}
In the official answer, \textbf{Travellers} is the large set, and the two means of travel, \textbf{train} (train travellers) and \textbf{bus} (bus travellers), are two separate groups inside it. They are shown as two disjoint circles, since a journey is made by one or the other, within one big circle. This is diagram (C).
\begin{center}
\begin{tikzpicture}[scale=0.7,font=\scriptsize]
\draw[fill=brand!8] (0,0) circle(1.6); \node at (0,1.25) {Travellers};
\draw[fill=white] (-0.7,-0.1) circle(0.6) node {Train}; \draw[fill=white] (0.7,-0.1) circle(0.6) node {Bus};
\end{tikzpicture}
\end{center}
\end{explain}
\end{question}

\begin{question}{46}
{How an achievement test is different from Entrance test ?\par}
\optsone{Achievement test needs high content validity whereas Entrance test needs high construct validity.}{Achievement test needs high construct validity whereas Entrance test needs high content validity.}{Achievement test needs high predictive validity whereas Entrance test needs high content validity.}{Achievement test needs high content validity whereas Entrance test needs high predictive validity.}
\answer{D}{Achievement test needs high content validity whereas Entrance test needs high predictive validity.}
\begin{explain}
An \textbf{achievement test} measures how well the content taught has been learnt, so it must have high \textbf{content validity} (it must sample the syllabus adequately). An \textbf{entrance test} is used to predict future success in a course, so it must have high \textbf{predictive validity}.
\end{explain}
\end{question}

\begin{question}{47}
{Sequential use of the three levels of coding - open coding, selective coding and theoretical coding is done in\par}
\optstwo{Phenomenological research}{Ethnographic research}{Grounded theory research}{Case study research}
\answer{C}{Grounded theory research}
\begin{explain}
The sequence \emph{open coding $\rightarrow$ selective coding $\rightarrow$ theoretical coding} belongs to Glaser's version of \textbf{grounded theory}. Strauss and Corbin's version uses open $\rightarrow$ axial $\rightarrow$ selective coding. In both, theory is generated from data by constant comparison.
\end{explain}
\end{question}

\begin{question}{48}
{In a reaction scale, the scale values are from 0 to 5. Which of the following scales of measurement is applicable for the data collected by this scale ?\par}
\optsfour{Nominal scale}{Ordinal scale}{Interval scale}{Ratio scale}
\answer{C}{Interval scale}
\begin{explain}
A reaction (rating) scale scored 0--5 has equal intervals between successive points. In educational research such summated rating data are treated as \textbf{interval-scale} data, which allows means and standard deviations to be computed. The official key accepts this answer.
\tip{Scales: Nominal (names or categories), Ordinal (rank order), Interval (equal intervals, arbitrary zero), Ratio (equal intervals with a true zero).}
\end{explain}
\end{question}

\begin{question}{49}
{Which of the following is a normal group ?\par}
\optsone{Skewness is zero and it is Platykurtic}{Mean = Median = Mode and it is Leptokurtic}{Mean = Median = Mode and it is Platykurtic}{Skewness is zero and it is Mesokurtic}
\answer{D}{Skewness is zero and it is Mesokurtic}
\begin{explain}
A \textbf{normal distribution} is perfectly symmetrical (\textbf{skewness $=0$}, mean $=$ median $=$ mode) and has \textbf{mesokurtic} peakedness (kurtosis $=3$, or excess kurtosis $=0$).
\why{Leptokurtic curves are more peaked and platykurtic curves flatter than normal.}
\end{explain}
\end{question}

\begin{question}{50}
{The transformation of human mind from theological metaphysical to a rational and scientific belief led to the emergence of\par}
\optstwo{Holism}{Rationalism}{Critical consciousness}{Positivism}
\answer{D}{Positivism}
\begin{explain}
Auguste Comte's \emph{Law of Three Stages} holds that human thought passes from the \textbf{theological} to the \textbf{metaphysical} and then to the \textbf{positive (scientific)} stage. This transformation to rational, scientific thinking gave rise to \textbf{Positivism}, the philosophy that only knowledge verified through observation and the scientific method is valid.
\end{explain}
\end{question}
